This Discussion interprets the results after they became available; the prespecified reading rules are archived with the code (\esmresource, Section~\ref{esm-sec:rule-record}). The central issue is how far discrimination between the simulated laws supports a claim about social influence. Three design limitations govern that interpretation: the restricted persistence alternatives, privileged knowledge of the response channel, and observation before the network process has had a burn-in.

\subsection{What classifier separation establishes}\label{sec:disc-meaning}

The confirmatory curve, $a^*(r)=\resultnum{j.b.r055.first.point}{0.10}$ for $r=0.55$--$0.90$ and \resultnum{j.b.r050.first.point}{0.30} at $r=0.50$, describes one classifier on three pooled long-lag correlations against trait-plus-autoregressive persistence. It holds with $r$ known, persistence matched to population moments, a correct response channel and panels as large as the SHP. AUC measures how well the classifier ranks panels drawn from these two simulated laws; the $0.70$ target is neither classification accuracy nor the power to reject persistence in one panel. Under the equal-variance normal approximation at a one-sided 5\% level, that target corresponds to single-panel power of about \resultnum{j.auc_translation.target070.power}{0.18} against a known persistence score distribution (Table~\ref{tab:auc-translation}). The curve therefore demonstrates conditional discrimination, without identifying social influence.

The parameter $a$ is a nominal innovation-variance fraction, not a realized share of attitude change or an empirical measure of social influence intensity. In the exploratory variance analysis, the influence term explains \resultnum{j.x2a.a010}{0.054} of the variance of the next latent attitude left after linear persistence at $a=0.10$, and \resultnum{j.x2a.a003}{0.016} at $a=0.03$. These are per-step latent quantities, not observed-scale effects. At the first-boundary cells, whose AUCs exceed $0.70$ on the coarse grid, even the oracle test rejected B for only \resultnum{j.x4a.boundary_rejection.min}{0.171} to \resultnum{j.x4a.boundary_rejection.max}{0.368} of single panels at $r\ge0.55$. It rejected B in at least four of five panels only from $a=0.30$ at $r\ge0.75$ (\resultnum{v4x.x4.x4a.r075.a3000.x4a.power_or}{0.826} to \resultnum{v4x.x4.x4a.r090.a3000.x4a.power_or}{0.944}), where influence explains \resultnum{j.x2a.a030}{0.167} of that residual variance; at $r=0.70$ and $a=0.30$ it rejected \resultnum{v4x.x4.x4a.r070.a3000.x4a.power_or}{0.759}. These exploratory rejection rates are references with privileged information, not the performance of an available empirical test.

The numerical boundary also depends on the target and the summaries. At the secondary target of \resultnum{v4.protocol.secondary_target.lower}{0.60} it falls to $a=\resultnum{v4.tab.T1sec.B.r075.a060.1}{0.03}$ at $r=0.75$ to $0.90$. In the exploratory comparison, thirty-six features reached the primary target at $a=\resultnum{j.x7.first.min}{0.003}$ to \resultnum{j.x7.first.max}{0.03} on the same panels on which three needed \resultnum{j.x7.frozen_first.min}{0.10} to \resultnum{j.x7.frozen_first.max}{0.30}. B was matched to none of the added features, and their gain changes estimation and regularization as well as information. Its source and specificity to influence were not isolated. The three-feature boundary is thus a property of the declared procedure, not a limit on all information in these panels.

\subsection{Restricted persistence alternatives}\label{sec:disc-mechanism}

The persistence alternatives do not span the temporal processes that could rival influence. Apart from transport and discretization, A, B and E share a plateau plus a single exponential decay, with no plateau for A. Their displayed boundaries cannot serve as independent corroboration, and the prespecified comparison across all persistence models was unavailable because C was invalid. D's boundaries were censored by matching gaps, so that comparison did not establish whether the classifier could distinguish influence from matched non-Gaussian state-dependent persistence.

The exploratory mean-field twin helps explain what this restricted family may miss. At $a=0.10$ the influence arm and its twin differed by at most \resultnum{j.x1a.lag_difference.a010.max}{0.0056} in mean lag correlation, yet the classifier separated them with AUCs of \resultnum{j.x1a.auc.a010.min}{0.623} to \resultnum{j.x1a.auc.a010.max}{0.716}. The twin keeps the effect through the own state while replacing the input from fixed contacts by fresh Gaussian noise. This contrast also changes dependence between people and the term's distribution, so it does not isolate a network-specific contribution. It is consistent with persistent contact input being an important source of separation, but does not establish that explanation (Section~\ref{sec:methods-identification}).

A persistent input can slow the decay of within-person correlations relative to a single exponential; B can mimic this with its plateau while missing the longer lags. In the free-share records, the influence arm's mean lag correlations at every boundary cell fell below those of known-share B at lags three to five, increasingly with the lag (at $r=0.70$, $\resultnum{j.x3.lag_shortfall.r070.l3}{-0.0011}$, $\resultnum{j.x3.lag_shortfall.r070.l4}{-0.0035}$ and $\resultnum{j.x3.lag_shortfall.r070.l5}{-0.0070}$). This is an observed lag pattern in the exploratory panels, not a decomposition of its causes.

No rival with serially correlated innovations, a second-order autoregression, or a sum of fast and slow autoregressive components was tested \cite{granger1976,asparouhov2018}. Nor did the design evaluate person-specific trends or persistence speeds \cite{bollencurran2004,asparouhov2018}, changes over the life course \cite{alwinkrosnick1991,lersch2023}, or period shocks shared by groups, which are relevant to trust in government. Such processes could reproduce part of the long-lag pattern and raise the boundary, possibly beyond the grid, but that possibility was not evaluated. Neither the mean-field twin nor the free-share analysis demonstrates what a matched second-order or two-component autoregressive rival would do. Testing those rivals, with heterogeneous persistence and matching at lags one to three, is needed before interpreting B's boundary as robust to richer individual dynamics.

\subsection{Privileged knowledge of response reliability}\label{sec:disc-attribution}

The confirmatory comparison gives both laws the same known signal share. With one answer per wave, estimating that share requires assumptions about the latent dynamics, so the restriction helps separate the laws through the very lag structure at issue (Section~\ref{sec:methods-identification}). A correctly specified channel shared by both arms changed little: \resultnum{j.sensitivity.response_channel.n_same_region}{44} of \resultnum{v4.inv.t4.n_valid}{45} sensitivity units kept the boundary interval above $a=0.03$. The exploratory analyses nevertheless show that allowing persistence to fit the channel, or misspecifying it, materially changes separation.

When B fitted its signal share to match lag three as well, most separation on lags four and five disappeared. The share moved by only \resultnum{j.x3.free.r_shift.min}{0.0011} to \resultnum{j.x3.free.r_shift.max}{0.0028} below the true value, more than halving B's trait weight. In the other direction, raising the share by \resultnum{j.x3.plus.offset}{0.01} moved that weight from \resultnum{j.x3.known.lambda.min}{0.139}--\resultnum{j.x3.known.lambda.max}{0.363} to \resultnum{j.x3.plus.lambda.min}{0.376}--\resultnum{j.x3.plus.lambda.max}{0.596}, with lags one and two matched in both cases. No free-share point AUC reached $0.70$ wherever that fit succeeded, although at $a=0.30$ and $r=0.90$ its interval reached $0.70$; in the four successful cells at $a=0.60$, the fit reduced B to an autoregression without trait. The share shift and extra matched lag are one fit whose contributions were not disentangled, and matching uncertainty was omitted. This is a demonstrated sensitivity of the lag-four-and-five classifier, not a revised confirmatory boundary or proof that the full laws are indistinguishable.

For the continuous response, a slightly lower share is equivalent to the true share with a short-lived latent component of corresponding size. Quasi-simplex models count such a component as error \cite{saris2014}, so this fit cannot tell slightly misjudged reliability from additional short-lived true variation. A share set \resultnum{j.x3.plus.offset}{0.01} too high put B's lag-three correlation \resultnum{j.x3.plus.l3_miss.min}{0.0069} to \resultnum{j.x3.plus.l3_miss.max}{0.0081} above the influence panels' at the boundary cells, the sign observed under influence. The response channel can also create separation when influence is absent: a misspecified channel reached AUCs up to \resultnum{j.x5.trend008.max}{0.907} (Table~\ref{tab:exploratory}). Unreliable items can make stable attitudes look changeable \cite{achen1975,ansolabehere2008}; a misspecified channel can also make persistence depart from the assumed model in influence's direction.

In that stress test, the arm without influence fell below B at lags three to five in all \resultnum{j.x5.signature.cells}{36} estimated cells, with a shortfall growing with the lag. Its median lag-five to lag-three ratio was \resultnum{j.x5.signature.ratio_l5_l3.median.mnar-ln15}{4.0} to \resultnum{j.x5.signature.ratio_l5_l3.median.trend-008}{4.8} by condition, against \resultnum{j.x3.lag_shortfall.ratio_l5_l3.min}{5.8} to \resultnum{j.x3.lag_shortfall.ratio_l5_l3.max}{6.9} under influence; individual cells overlap (\esmresource, Table~\ref{esm-tab:x5-lags}). Whether these departures would fool a score trained on influence was not evaluated. Neither the stress test nor the free-share contrast was run with the thirty-six features. They explain fragility within their declared comparisons, without establishing that every classifier would have the same sensitivity.

The SHP trust item illustrates the difficulty of supplying the privileged information (\esmresource, Section~\ref{esm-sec:shp}). An eleven-wave quasi-simplex model estimates reliability in \resultnum{v4.protocol.shp.anchor_year}{2004} at \resultnum{v4.shp.estimate}{0.674} and a trait--state--error fit at \resultnum{j.shp.starts.estimate}{0.696}, but the quasi-simplex fit is strained (persistence \resultnum{j.shp.multiwave.persistence}{0.961}; residual correlations up to \resultnum{j.shp.multiwave.max_cc_residual}{0.078}). Both models count stable response style as signal, whereas the simulation treats it as error: they estimate reliability, while the simulated signal share is closer to measurement quality \cite{saris2014}. If the latent process is a first-order autoregression with independent errors, Heise's coefficient \cite{heise1969} estimates a three-wave window's middle-wave reliability. Eight of the nine windows of \resultnum{j.shp.heise.w1999.start}{1999}--\resultnum{j.shp.heise.w2007.end}{2009} lie within \resultnum{j.shp.heise.others.spread}{0.029} of each other (standard errors \resultnum{j.shp.heise.w2004.se}{0.011} to \resultnum{j.shp.heise.w2002.se}{0.014}). The window of \resultnum{j.shp.heise.w2003.start}{2003}--\resultnum{j.shp.heise.w2003.end}{2005}, centered on refreshment-sample entry, lies \resultnum{j.shp.heise.dip2003}{0.065} below its neighbors' mean, giving a spread of \resultnum{j.shp.heise.spread}{0.077} including it. Sampling error, composition and departures from the model may all contribute; the spread does not estimate drift.

These standard errors provide a comparison of scales with the simulation's third-decimal sensitivity, not an empirical reliability threshold. A trend from $r-0.03$ to $r+0.03$ already separated simulated panels without influence. The short-lived response component, including error and short-lived true variation, is identified from one answer per wave only through the lag structure the classifier reads; assessing it at that precision would require several indicators or repeated measurement within a wave. Under this generator and channel, shares near the item's estimates gave $a=\resultnum{j.b.r065.first.point}{0.10}$ against B. This is not an estimated threshold for the SHP item: its lag profile, heaping at 0, 5 and 10, and attrition were never compared with the synthetic panels. An empirical analyst must also estimate persistence from one panel, rather than population moments. No valid analyst-implementable single-panel procedure was evaluated, so the curve does not establish the power available in that setting.

\subsection{Observation without network burn-in}\label{sec:disc-initialization}

The influence process begins from states unrelated to the network. With per-wave persistence of 0.85 and contacts fixed from wave 0, dependence between a person and those contacts plausibly takes several waves to build. The observed waves include that adjustment, which could contribute to the lag pattern used for discrimination. No analysis compared these panels with panels observed after burn-in. The contribution of the transient, and the direction in which burn-in would move the boundary, are therefore unknown. The mean-field contrast does not isolate this contribution either.

Other generator assumptions also limit transfer to empirical panels: one persistence and repulsion setting, eight contacts drawn uniformly at random and fixed for eleven years, perfect perception of their states, one step per wave, no coupling within households, and attrition that is neither selective nor temporary. The influence arm has no stable trait, whereas the SHP anchor's trait--state--error fit puts \resultnum{j.shp.starts.trait_share_2004}{0.758} of latent variance in a trait (\esmresource, Section~\ref{esm-sec:shp-results}). Influence confined to a smaller state component could require stronger input; fewer contacts could enlarge the persistent input and lower the boundary; homophily \cite{mcpherson2001}, clustering, contact turnover and noisy perception could fold it into own persistence or fresh noise and raise the boundary. These are untested mechanism-based conjectures, not sensitivity results. Adding a stable trait and comparing initialization with and without burn-in would address distinct parts of this scope.

\subsection{Relation to earlier work}\label{sec:disc-literature}

The ambiguity that Kiley and Vaisey describe between frequent updating and settled dispositions measured with error \cite{kiley2020} also arises between network influence and persistence, and turns on the response channel. Comparisons with no-change baselines \cite{moor2026,coding2024} do not match persistence to the same short-term stability; the boundary is a trajectory-level counterpart of the minimum change that Moor-Smith and Carpentras find a generator needs to beat no change. The stress test is a dynamic counterpart of measurement distortions in opinion-dynamics models \cite{carpentras2022,carpentras2023} and of spurious cross-lagged effects without a stable trait \cite{lucas2023}.

The result also raises a question about calibrations such as Gestefeld and Lorenz's to the same SHP item, in which each agent interacts with a randomly selected agent, or returns to its initial opinion, about 24 times per survey year \cite{gestefeld2023}. No persistent input from the same contacts builds up. At their best fits' return probability of 0.1 per step, about nine in ten agents return between annual waves, so year-to-year dependence runs largely through the initial opinion, a stable trait in panel terms. The closest analog here is the mean-field twin, which redraws the contact input at every step; against B it reached the target only at $r=0.90$, at $a=\resultnum{j.x1b.r090.first.point}{0.30}$. Neither the Gaussian twin nor the 24-step or anchoring sensitivity represents their model. If persistent input carries much of the present separation, random encounters could leave less separable panel transitions; this is a hypothesis for those calibrations, not a result about them. Their transition fit alone does not show that interaction is needed without comparison with matched persistence passed through the same response channel (Section~\ref{sec:implications}).

Without exposure data or other identifying restrictions, panels like these can establish a departure only from the persistence models considered. This meets the older problem of distinguishing contagion or state dependence from heterogeneity \cite{feller1943,heckman1981,gelastopoulos2025}; even with observed networks, separating influence from correlated effects and homophily requires strong assumptions \cite{manski1993,shalizi2011}.

\subsection{Further scope, uncertainty and reproducibility limits}\label{sec:limitations}

With 10,626 people and eleven waves, pooled correlations are very precise. Fewer people or waves would be expected to reduce that precision, and three-wave panels have no lags three to five. The grid starts at $r=0.50$, while many single items are less reliable \cite{hout2016,ansolabehere2008}. The channel is classical error around a latent state, discretized by within-wave ranks; answers that sample momentary considerations \cite{zaller1992} were not modeled, and the attrition result depends on pooling all available pairs of answers.

Inference at positive strength rests on limited checks. Band coverage was checked over the grid only at $r=0.70$, and pointwise at $a=0.03$ with $r=0.50$ and $0.90$. No undercoverage was detected for complete gap-free bands or at those pointwise cells; counting gaps as misses failed the coverage criterion (Section~\ref{sec:results-exploratory}). Near $a=0.03$, two gate diagnostics point in opposite directions (\esmresource, Section~\ref{esm-sec:record-gate}). With point estimates and critical values held fixed, an upper band at $a=0.03$ would reach $0.70$ only with a standard error about \resultnum{j.b.margin.a0300.se_factor.min}{2.0} times larger (at $r=\resultnum{j.b.margin.a0300.se_factor.min.r}{0.90}$; more elsewhere), whereas the nuisance-variance check at $r=0.70$ implies growth of at most about \resultnum{j.gate.g2c.a0300.se_growth_max}{1.46} times, and its applicability at other shares was not assessed. The zero-width interval at $r=0.65$ could widen by one grid step within the dispersion observed at $a=0.10$ and $r=0.70$.

The evaluation was not blind: at $r=0.70$ an earlier calibration proposal had shown an AUC of \resultnum{j.exposure.proposal_a.a010.mean}{0.756} at $a=0.10$ and \resultnum{j.exposure.proposal_a.a003}{0.589} at $a=0.03$, anticipating the boundary at that share. The AUC target and $0.03$ cut preceded these probes but not all earlier calibration (\esmresource, Section~\ref{esm-sec:record-exposure}). New seeds limit adaptation without removing it; the new confirmatory content is mainly the curve at the other eight shares.

Execution is bound to one recorded computing environment. Components agreed with independent implementations, but replay of the full pipeline has not been independently demonstrated, three frozen tests fail on Linux, and after an operating-system update on 9 October, that machine itself no longer passes the freeze check. The library entry points provide a route to rerun the simulation outside that environment, with possible floating-point and tolerance-edge differences. Agreement of such a complete rerun with the released results remains to be assessed (Data and code availability; \esmresource, Section~\ref{esm-sec:verification}).

\subsection{Implications for empirical work}\label{sec:implications}

For calibrating opinion-dynamics models to panel surveys, the findings support comparing influence with persistence matched to the same short-term stability, judging fit on patterns not used for calibration \cite{grimm2005}, and simulating the response channel, including plausible misspecifications, in every model. Letting persistence fit its signal share tests a restriction that the confirmatory comparison held fixed. Reporting AUC alongside an oracle single-panel rejection rate can expose the difference between discrimination of laws and power in one panel, while retaining the oracle rate's privileged-information caveat.

The archived code can be adapted outside its frozen part by replacing this study's generator with the calibrated model. Features should include those central to the empirical claim, such as one-wave transitions; the frozen channel's equiprobable categories discard the item's marginal distribution, and lag correlations should use answers ranked or standardized within waves. Any claim of social influence also needs item-specific evidence on reliability, its stability across waves and item-related attrition \cite{lipps2007,voorpostel2010}, together with restrictions that distinguish social input from persistent nonsocial inputs such as a media diet or local conditions. In the SHP, co-resident pairs provide a measured relationship that actor--partner models can use \cite{voorpostel2025shp,kennykashy2006}. Such evidence addresses attribution beyond the classifier comparison performed here.
